\documentclass[12pt]{article}
\usepackage[T1]{fontenc}
\usepackage[utf8]{inputenc}
\usepackage{lmodern}
\usepackage{textcomp}
\usepackage{enumitem}
\usepackage{geometry}
\geometry{margin=1in}
\begin{document}
\begin{center}
Answer Key - Sociology - Graduate Level Assessment 11 \
\small (Answers are ordered exactly as in the question paper.)
\end{center}

\section*{Multiple Choice}
\begin{enumerate}[label=\arabic*.]
\item \textbf{Answer:} D
\\ \textit{Options:} A) this pattern of speech made them the target of bullying; B) they referred to explicit, context independent meanings; C) they prevented children from communicating outside of their peer groups; D) they involved short, simple sentences with a small vocabulary
\\ \textit{Note:} Correct option: D - they involved short, simple sentences with a small vocabulary
\item \textbf{Answer:} D
\\ \textit{Options:} A) flexibility was reduced by the introduction of detailed daily worksheets; B) the state had greater control than managers over production processes; C) ownership was transferred from small shareholders to senior managers; D) employment depended on performance rather than trust and commitment
\\ \textit{Note:} Correct option: D - employment depended on performance rather than trust and commitment
\item \textbf{Answer:} D
\\ \textit{Options:} A) it is committed on a larger, often global, scale, and is well organized; B) it is associated with political conflict between states and their citizens; C) it can have far-reaching effects upon international relations; D) all of the above
\\ \textit{Note:} Correct option: D - all of the above
\item \textbf{Answer:} D
\\ \textit{Options:} A) an urban set, involved in civic bodies and voluntary associations; B) too diverse to have a strong sense of class consciousness; C) often involved in 'white collar' work; D) all of the above
\\ \textit{Note:} Correct option: D - all of the above
\item \textbf{Answer:} C
\\ \textit{Options:} A) participated in the same leisure pursuits and events of the 'social calendar'; B) emulated the lifestyle and cultural values of the traditional aristocracy; C) owned companies and financial assets that generated wealth through corporations; D) had direct, personal ownerships of land and businesses as physical assets
\\ \textit{Note:} Correct option: C - owned companies and financial assets that generated wealth through corporations
\end{enumerate}

\section*{True / False}
\begin{enumerate}[label=\arabic*.]
\setcounter{enumi}{5}
\item \textbf{Answer:} False
\\ \textit{Options:} A) True; B) False
\\ \textit{Note:} LLM-generated false statement. Correct answer: 'chronic, degenerative diseases, such as cancer, heart disease, and strokes'.
\item \textbf{Answer:} True
\\ \textit{Options:} A) True; B) False
\\ \textit{Note:} LLM-generated true statement based on MCQ.
\item \textbf{Answer:} False
\\ \textit{Options:} A) True; B) False
\\ \textit{Note:} LLM-generated false statement. Correct answer: 'universalistic benefits and public sector employment'.
\item \textbf{Answer:} True
\\ \textit{Options:} A) True; B) False
\\ \textit{Note:} LLM-generated true statement based on MCQ.
\item \textbf{Answer:} True
\\ \textit{Options:} A) True; B) False
\\ \textit{Note:} LLM-generated true statement based on MCQ.
\end{enumerate}

\section*{Long Form Response}
\begin{enumerate}[label=\arabic*.]
\setcounter{enumi}{10}
\item \textbf{Answer:} the extended power of nation states. Explain why this is the correct answer and provide supporting evidence. Reference key facts or concepts that support this choice.
\\ \textit{Note:} Long-form derived from MCQ; award credit for stating the correct idea and giving supporting reasoning.
\item \textbf{Answer:} a unit with a division of labour that extends across ethnic and cultural groups. Explain why this is the correct answer and provide supporting evidence. Reference key facts or concepts that support this choice.
\\ \textit{Note:} Long-form derived from MCQ; award credit for stating the correct idea and giving supporting reasoning.
\end{enumerate}

\vfill
\noindent\textit{End of Answer Key}
\end{document}